% GENERATED FILE. Edit canonical lessons, then run npm run generate:books.
% canonical-source-hash: fnv1a64:c6d16943c3b8285e
% canonical-lessons: ZH-C525-zhoudao, ZH-C525-gewai, ZH-W525-ge-pattern, ZH-C525-shouhuoyuan, ZH-C525-tousu, ZH-C525-tuikuan, ZH-R525-review-520-525

\chapter{At Your Service}
\label{ch:zh-at-your-service}
\hlchaptermodality{525}

\begin{chapteropening}
\textbf{By the end of this chapter:} \emph{I can talk about attentive service and shop assistants, complain and get a refund, and say something is especially good.}

Complete the last lesson of chapter 525: I can talk about attentive service and shop assistants, complain and get a refund, and say something is especially good.
\end{chapteropening}

\section[zhōudào]{zhōudào — attentive and thorough (of service)}

\label{lesson:ZH-C525-zhoudao}



\begin{quote}
Before the new one: say the Mandarin for to offer a lower price back, then the Mandarin for the cost of living.
\end{quote}

\subsection*{You'll want to know: zhōudào}
\emph{zhōudào} — \textquotedblleft{}attentive and thorough (of service)\textquotedblright{}.

\begin{quote}
\emph{zhè jiā fàndiàn de fúwùyuán bǐ nà jiā zhōudào} — The waiters in this restaurant are more attentive than in that one.
\end{quote}

Its characters come later; for now it is a word for the ear.

\begin{cousinweb}
\emph{zhōudào} is built from \emph{zhōu} and \zh{到} 'to reach', both inherited from Middle Chinese, \emph{tsyuw} and \emph{tawH} (Guangyun, in Baxter's transcription), which glosses the first 'all round', 'complete' and 'close' and the second 'to arrive': reaching all round.
\end{cousinweb}

\begin{grammarlens}[title={What you've built}]
One more word for saying how it was.
\end{grammarlens}

\subsection*{Your turn}
\begin{itemize}
\raggedright
  \item \emph{Say it:} \emph{zhōudào}
  \item \emph{Say it:} \emph{zhōudào}, once more
  \item \emph{Say it:} \emph{zhè jiā fàndiàn de fúwùyuán bǐ nà jiā zhōudào} — The waiters in this restaurant are more attentive than in that one.
  \item \emph{Recall:} say the Mandarin for to offer a lower price back, then the Mandarin for the cost of living, then say \emph{zhōudào} again
\end{itemize}

\begin{tcolorbox}[breakable,colback=teal!4,colframe=teal!35!black,title={Before you move on}]
What does zhōudào mean? (\textquotedblleft{}Attentive and thorough (of service)\textquotedblright{}.) Say it once more.
\end{tcolorbox}
\section[\texorpdfstring{\zh{格外} (géwài)}{格外 (géwài)}]{\zh{格外} (géwài) — more than usual; unusually}

\label{lesson:ZH-C525-gewai}



\begin{quote}
Before the new one: say the Mandarin for the cost of living, then the Mandarin for attentive and thorough.
\end{quote}

\subsection*{You'll want to know: \zh{格外}}
\textbf{\zh{格外}} — \emph{géwài} — \textquotedblleft{}more than usual; unusually\textquotedblright{}.

\begin{quote}
\textbf{\zh{今天的菜格外好吃}} \emph{jīntiān de cài géwài hǎochī} — The food is especially good today.
\end{quote}

The new character, in pieces: \zh{格} — ten strokes in three pieces: on the left the tree piece, four strokes, as in \zh{机}; on the right the walking foot, three strokes, as in the middle of \zh{客}; and \zh{口} under it, three strokes. The next lesson puts it in your hand.

\begin{cousinweb}
\textbf{\zh{格外}} (\emph{géwài}) is built from \zh{格} and \zh{外} 'outside', both inherited from Middle Chinese, \emph{kaek} and \emph{ngwajH} (Guangyun, in Baxter's transcription). The dictionary glosses the first 'a pattern' and 'a measure', and lists another reading, \emph{kak}, for a branch; outside the measure is beyond the usual, though how the pair came to mean 'especially' is not settled.
\end{cousinweb}

\begin{grammarlens}[title={What you've built}]
One more word for saying how it was.
\end{grammarlens}

\subsection*{Your turn}
\begin{itemize}
\raggedright
  \item \emph{Say it:} \emph{géwài}
  \item \emph{Say it:} \emph{géwài}, once more
  \item \emph{Say it:} \emph{jīntiān de cài géwài hǎochī} — The food is especially good today.
  \item \emph{Recall:} say the Mandarin for the cost of living, then the Mandarin for attentive and thorough, then say \emph{géwài} again
\end{itemize}

\begin{tcolorbox}[breakable,colback=teal!4,colframe=teal!35!black,title={Before you move on}]
What does \zh{格外} mean? (\textquotedblleft{}More than usual; unusually\textquotedblright{}.) Say it once more.
\end{tcolorbox}
\section[\texorpdfstring{\zh{格} (gé)}{格 (gé)}]{\zh{格} — the tree piece, then the walking foot over \zh{口}}

\label{lesson:ZH-W525-ge-pattern}



\begin{quote}
The word first, then the character inside it.

\begin{itemize}
\raggedright
  \item \emph{Recall:} say \emph{géwài} — \textbf{R1}, one lesson back, before it decays
\end{itemize}

\begin{quote}
\zh{客} \zh{机}
\end{quote}

\emph{Read it:} each one aloud, with nothing said first — \emph{kè}, \emph{jī}
\end{quote}

\subsection*{Script — the tree piece, then the walking foot over \zh{口}}
\hlblockfigure{\detokenize{figures/ZH-W525-ge-pattern-filmstrip.pdf}}{How \zh{格} is written, stroke by stroke}

\begin{quote}
\zh{格}
\end{quote}

\textbf{Three pieces: the tree piece on the left, then the walking foot, then \zh{口} under it.} Count the pieces, not the strokes.

\textbf{The left piece, the tree piece, four strokes}, narrow, as in \zh{机}: a short bar, a vertical down through it, a stroke falling to the left, then a short dot to the right.

\textbf{The walking foot, three strokes}, as in \zh{客}: a short stroke falling to the left, a bar that turns into a long stroke falling to the left, then a long stroke falling to the right.

\textbf{\zh{口}, three strokes}, at the foot: a vertical, a bar that turns down, and a bar closing it.

4 + 3 + 3 = ten.

\subsection*{Your turn}
\begin{itemize}
\raggedright
  \item \emph{Write it:} \zh{格}, all ten, with the model in front of you
  \item \emph{Cover:} write it again from memory, and count the pieces, not the strokes
\end{itemize}

\begin{tcolorbox}[breakable,colback=teal!4,colframe=teal!35!black,title={Before you move on}]
How many pieces does \zh{格} have? (\textbf{Three: the tree piece, the walking foot and \zh{口}}.) Which character has the same right side at its foot? (\textbf{\zh{客}}.)

Source: \href{https://raw.githubusercontent.com/chanind/hanzi-writer-data/68d10a4b21150cae5e1ebbd223eed289cf32d90c/data/\%E6\%A0\%BC.json}{Hanzi Writer Data \zh{格}, pinned stroke paths and medians}.
\end{tcolorbox}
\section[shòuhuòyuán]{shòuhuòyuán — shop assistant}

\label{lesson:ZH-C525-shouhuoyuan}



\begin{quote}
Before the new one: say the Mandarin for attentive and thorough, then the Mandarin for more than usual.
\end{quote}

\subsection*{You'll want to know: shòuhuòyuán}
\emph{shòuhuòyuán} — \textquotedblleft{}shop assistant\textquotedblright{}.

\begin{quote}
\emph{wǒ juéde zhèr de shòuhuòyuán zuì rèqíng} — I think the shop assistants here are the friendliest.
\end{quote}

Its characters come later; for now it is a word for the ear.

\begin{cousinweb}
\emph{shòuhuòyuán} is built from \emph{shòu} 'to sell', \emph{huò} 'goods' and \emph{yuán} 'a member of staff', all inherited from Middle Chinese, \emph{dzyuwH}, \emph{xwaH} and \emph{hjwen} (Guangyun, in Baxter's transcription). The dictionary glosses the first 'to sell goods', the second 'wealth' and the third as the number of things; the sense 'staff' is not given there.
\end{cousinweb}

\begin{grammarlens}[title={What you've built}]
One more word for saying how it was.
\end{grammarlens}

\subsection*{Your turn}
\begin{itemize}
\raggedright
  \item \emph{Say it:} \emph{shòuhuòyuán}
  \item \emph{Say it:} \emph{shòuhuòyuán}, once more
  \item \emph{Say it:} \emph{wǒ juéde zhèr de shòuhuòyuán zuì rèqíng} — I think the shop assistants here are the friendliest.
  \item \emph{Recall:} say the Mandarin for attentive and thorough, then the Mandarin for more than usual, then say \emph{shòuhuòyuán} again
\end{itemize}

\begin{tcolorbox}[breakable,colback=teal!4,colframe=teal!35!black,title={Before you move on}]
What does shòuhuòyuán mean? (\textquotedblleft{}Shop assistant\textquotedblright{}.) Say it once more.
\end{tcolorbox}
\section[tóusù]{tóusù — to make a formal complaint}

\label{lesson:ZH-C525-tousu}



\begin{quote}
Before the new one: say the Mandarin for more than usual, then the Mandarin for a shop assistant.
\end{quote}

\subsection*{You'll want to know: tóusù}
\emph{tóusù} — \textquotedblleft{}to make a formal complaint\textquotedblright{}.

\begin{quote}
\emph{wǒ tóusù guo zhè jiā shāngdiàn} — I've complained about this shop before.
\end{quote}

Its characters come later; for now it is a word for the ear.

\begin{cousinweb}
\emph{tóusù} is built from \emph{tóu} 'to throw, to hand in' and \emph{sù} 'to accuse', both inherited from Middle Chinese, \emph{duw} and \emph{suH} (Guangyun, in Baxter's transcription), which glosses the first 'to entrust' and 'to throw' and the second 'to bring a case' and 'to tell'.
\end{cousinweb}

\begin{grammarlens}[title={What you've built}]
One more word for saying how it was.
\end{grammarlens}

\subsection*{Your turn}
\begin{itemize}
\raggedright
  \item \emph{Say it:} \emph{tóusù}
  \item \emph{Say it:} \emph{tóusù}, once more
  \item \emph{Say it:} \emph{wǒ tóusù guo zhè jiā shāngdiàn} — I've complained about this shop before.
  \item \emph{Recall:} say the Mandarin for more than usual, then the Mandarin for a shop assistant, then say \emph{tóusù} again
\end{itemize}

\begin{tcolorbox}[breakable,colback=teal!4,colframe=teal!35!black,title={Before you move on}]
What does tóusù mean? (\textquotedblleft{}To make a formal complaint\textquotedblright{}.) Say it once more.
\end{tcolorbox}
\section[tuìkuǎn]{tuìkuǎn — to give the money back; to refund}

\label{lesson:ZH-C525-tuikuan}



\begin{quote}
Before the new one: say the Mandarin for a shop assistant, then the Mandarin for to make a formal complaint.
\end{quote}

\subsection*{You'll want to know: tuìkuǎn}
\emph{tuìkuǎn} — \textquotedblleft{}to give the money back; to refund\textquotedblright{}.

\begin{quote}
\emph{wǒ děng le yí ge xīngqī tāmen cái tuìkuǎn} — I waited a week before they gave me my money back.
\end{quote}

Its characters come later; for now it is a word for the ear.

\begin{cousinweb}
\emph{tuìkuǎn} is built from \emph{tuì} 'to go back' and \emph{kuǎn}, both inherited from Middle Chinese, \emph{thwojH} and \emph{khwanX} (Guangyun, in Baxter's transcription). The dictionary glosses the first 'to retreat' and the second 'sincere' and 'to knock'; the sense 'a sum of money' is not given there.
\end{cousinweb}

\begin{grammarlens}[title={What you've built}]
One more word for saying how it was.
\end{grammarlens}

\subsection*{Your turn}
\begin{itemize}
\raggedright
  \item \emph{Say it:} \emph{tuìkuǎn}
  \item \emph{Say it:} \emph{tuìkuǎn}, once more
  \item \emph{Say it:} \emph{wǒ děng le yí ge xīngqī tāmen cái tuìkuǎn} — I waited a week before they gave me my money back.
  \item \emph{Recall:} say the Mandarin for a shop assistant, then the Mandarin for to make a formal complaint, then say \emph{tuìkuǎn} again
\end{itemize}

\begin{tcolorbox}[breakable,colback=teal!4,colframe=teal!35!black,title={Before you move on}]
What does tuìkuǎn mean? (\textquotedblleft{}To give the money back; to refund\textquotedblright{}.) Say it once more.
\end{tcolorbox}
\section[(dialogue)]{Looking back: from Around the House to At Your Service}

\label{lesson:ZH-R525-review-520-525}



\begin{quote}
Say the Mandarin for a household appliance and a light switch.
\end{quote}

\subsection*{Your turn}
\begin{itemize}
\raggedright
  \item \emph{Say it:} a household appliance, a light switch, a water heater, central heating, to use little electricity
  \item \emph{Say it:} to scrub clean, to sew by hand, to lay out a room, neat and orderly, hygiene
  \item \emph{Say it:} to leak water, to call at the house, the water bill, a doorbell, the floor of a room
  \item \emph{Say it:} to sell wholesale, buying and selling for a living, the seller (in a sale), the buyer (in a sale), to sell retail
  \item \emph{Say it:} to settle the bill, at half price, to name one's price, to offer a lower price back, the cost of living
  \item \emph{Say it:} attentive and thorough, more than usual, a shop assistant, to make a formal complaint, to give the money back
  \item \emph{Read it:} \zh{卖} \zh{价} \zh{格} — each one aloud, then a word that holds it
\end{itemize}

\begin{tcolorbox}[breakable,colback=teal!4,colframe=teal!35!black,title={Before you move on}]
A household appliance? (\textbf{\zh{家电}}, \emph{jiādiàn}.) A light switch? (\textbf{\zh{开关}}, \emph{kāiguān}.) A water heater? (\emph{rèshuǐqì}.) Central heating? (\emph{nuǎnqì}.) To use little electricity? (\emph{shěngdiàn}.) To scrub clean? (\emph{shuā}.) To sew by hand? (\emph{féng}.) To lay out a room? (\emph{bùzhì}.) Neat and orderly? (\emph{zhěngqí}.) Hygiene? (\emph{wèishēng}.) To leak water? (\emph{lòushuǐ}.) To call at the house? (\textbf{\zh{上门}}, \emph{shàngmén}.) The water bill? (\emph{shuǐfèi}.) A doorbell? (\emph{ménlíng}.) The floor of a room? (\emph{dìbǎn}.) To sell wholesale? (\emph{pīfā}.) Buying and selling for a living? (\textbf{\zh{买卖}}, \emph{mǎimai}.) The seller (in a sale)? (\textbf{\zh{卖家}}, \emph{màijiā}.) The buyer (in a sale)? (\textbf{\zh{买家}}, \emph{mǎijiā}.) To sell retail? (\emph{língshòu}.) To settle the bill? (\emph{jiézhàng}.) At half price? (\textbf{\zh{半价}}, \emph{bànjià}.) To name one's price? (\textbf{\zh{要价}}, \emph{yàojià}.) To offer a lower price back? (\textbf{\zh{还价}}, \emph{huánjià}.) The cost of living? (\emph{wùjià}.) Attentive and thorough? (\emph{zhōudào}.) More than usual? (\textbf{\zh{格外}}, \emph{géwài}.) A shop assistant? (\emph{shòuhuòyuán}.) To make a formal complaint? (\emph{tóusù}.) To give the money back? (\emph{tuìkuǎn}.)
\end{tcolorbox}
